\textbf{Время:} 01:15

\textbf{Тезис:} фемтоскопия пионов чувствительна к области их последнего взаимодействия и рождения при распадах.

\textbf{Визуал:} ранняя плотная стадия → расширение → разные точки и времена испускания.

\textbf{Текст:} Пион может неоднократно рассеяться и изменить импульс. Поэтому конечный импульс не является простой меткой того места, где он впервые появился.

По мере расширения реакции редеют. Химическое вымораживание связано с прекращением эффективного изменения выходов частиц, кинетическое — с прекращением перерассеяния. Химическое вымораживание не останавливает последующие распады резонансов. Конечный пион мог родиться при таком распаде.

Для разных частиц это происходит в разных точках и в разное время. Источник в нашем дальнейшем расчёте — именно такое распределение испускания, а не мгновенная фотография раннего столкновения. Когда сравниваем радиусы, важно помнить, какую стадию и какой набор частиц мы выбрали.

\textbf{Переход:} Наша область испускания уже исчезла. Какая информация о ней остаётся доступной?
